\documentclass[handout,aspectratio=169]{beamer}

%------------------------------------------------
% PAQUETES
%------------------------------------------------
\usepackage[spanish,es-noshorthands]{babel}
\usepackage[utf8]{inputenc}
\usepackage[T1]{fontenc}
\usepackage{lmodern}
\usepackage{amsmath, amssymb}
\usepackage{tikz}
\usepackage{xcolor}
\usepackage{tcolorbox}
\usepackage{booktabs}
\usepackage{array}
\usepackage{ragged2e}
\usepackage{multicol}

%------------------------------------------------
% TEMA
%------------------------------------------------
\usetheme{Madrid}
\useinnertheme{rounded}
\setbeamertemplate{navigation symbols}{}

%------------------------------------------------
% COLORES PSICOLOGÍA / UST
%------------------------------------------------
\definecolor{Primary}{HTML}{6A1B9A}
\definecolor{Secondary}{HTML}{8E24AA}
\definecolor{LightBg}{HTML}{F3E5F5}
\definecolor{DarkText}{HTML}{263238}
\definecolor{AlertRed}{HTML}{C62828}

\setbeamercolor{structure}{fg=Primary}
\setbeamercolor{frametitle}{bg=Primary,fg=white}
\setbeamercolor{title}{fg=white,bg=Primary}
\setbeamercolor{block title}{bg=Primary,fg=white}
\setbeamercolor{block body}{bg=LightBg,fg=black}
\setbeamercolor{alerted text}{fg=AlertRed}

%------------------------------------------------
% FOOTLINE
%------------------------------------------------
\setbeamertemplate{footline}{
	\leavevmode%
	\hbox{%
		\begin{beamercolorbox}[wd=.82\paperwidth,ht=2.5ex,dp=1ex,left]{author in head/foot}
			\hspace{1em} Pensamiento Matemático Ciencias Sociales | Psicología
		\end{beamercolorbox}%
		\begin{beamercolorbox}[wd=.18\paperwidth,ht=2.5ex,dp=1ex,center]{date in head/foot}
			\insertframenumber/\inserttotalframenumber
	\end{beamercolorbox}}%
}

%------------------------------------------------
% CAJAS
%------------------------------------------------
\newtcolorbox{idea}{
	colback=LightBg,
	colframe=Primary,
	arc=2mm,
	boxrule=0.8pt,
	title=\textbf{Idea clave}
}

\newtcolorbox{alerta}{
	colback=red!5,
	colframe=AlertRed,
	arc=2mm,
	boxrule=0.8pt,
	title=\textbf{Cuidado}
}

\newtcolorbox{actividad}{
	colback=blue!5,
	colframe=blue!60!black,
	arc=2mm,
	boxrule=0.8pt,
	title=\textbf{Actividad}
}

%------------------------------------------------
% DATOS
%------------------------------------------------
\title[Potencias y Álgebra]{Potencias, Expresiones Algebraicas y Patrones}
\subtitle{Aplicaciones contextualizadas a Psicología}
\author{Prof. Luis Tulio González Alcaino}
\institute{Universidad Santo Tomás \\ Departamento de Ciencias Básicas}
\date{\today}

%------------------------------------------------
\begin{document}
	%------------------------------------------------
	
	\begin{frame}
		\titlepage
	\end{frame}
	
	%------------------------------------------------
	
	\begin{frame}{Ruta de la clase}
		\tableofcontents
	\end{frame}
	
	%------------------------------------------------
	\section{Objetivos de aprendizaje}
	
	\begin{frame}{Objetivos de aprendizaje}
		Al finalizar la clase, el estudiante será capaz de:
		
		\begin{enumerate}
			\item Reconocer el concepto de potencia y sus propiedades.
			\item Identificar término algebraico, expresión algebraica y término semejante.
			\item Aplicar propiedades de potencias en situaciones contextualizadas a Psicología.
			\item Simplificar expresiones algebraicas asociadas a modelos psicológicos simples.
			\item Integrar patrones y relaciones matemáticas en escalas, registros y mediciones.
			\item Detectar inconsistencias frecuentes en procedimientos algebraicos.
		\end{enumerate}
	\end{frame}
	
	%------------------------------------------------
	\section{Concepto de potencia}
	
	\begin{frame}{Concepto de potencia}
		\begin{idea}
			Una potencia representa una multiplicación reiterada de un mismo factor.
		\end{idea}
		
		\[
		a^n=
		\underbrace{a\cdot a\cdot a\cdots a}_{n\text{ veces}}
		\]
		
		\pause
		
			\begin{itemize}
				\item \(a\): base.
				\item \(n\): exponente.
				\item \(a^n\): potencia.
			\end{itemize}
	
		\pause
		
		\begin{exampleblock}{Ejemplo en Psicología}
			Si una tarea de memoria duplica su dificultad en cada nivel, el número relativo de combinaciones puede modelarse como: $2^n$, donde \(n\) representa el nivel de dificultad.
		\end{exampleblock}
	\end{frame}
	
	%------------------------------------------------
	
	\begin{frame}{Ejemplo: carga cognitiva}
		En una prueba experimental, la cantidad de estímulos visuales aumenta duplicándose en cada etapa. $2^1=2,\quad 2^2=4,\quad 2^3=8,\quad 2^4=16$
		
		\pause
		
		\begin{block}{Interpretación}
			Cada nuevo nivel exige procesar el doble de estímulos que el nivel anterior.
		\end{block}
		
		\pause
		
		\begin{actividad}
			Complete la tabla:
			\[\begin{array}{c|c}
				\text{Nivel} & \text{Estímulos} \\
				\hline
				1 & 2 \\
				2 & 4 \\
				3 & 8 \\
				4 & 16 \\
				5 & ? \\
				6 & ?
			\end{array}
			\]
		\end{actividad}
	\end{frame}
	
	%------------------------------------------------
	\section{Propiedades de las potencias}
	
	\begin{frame}{Producto de potencias de igual base}
		\begin{idea}
			Cuando se multiplican potencias de igual base, se conserva la base y se suman los exponentes.
		\end{idea}
		
		\[
		a^m\cdot a^n=a^{m+n}
		\]
		
		\pause
		
		\begin{exampleblock}{Ejemplo contextualizado}
			En un modelo de registro conductual, una variable \(x\) representa la frecuencia base de una conducta observada. 
			$x^3\cdot x^4=x^{3+4}=x^7$			
		\end{exampleblock}
		
		\pause
		
		\begin{block}{Interpretación}
			Se combinan dos factores de crecimiento asociados a la misma variable psicológica.
		\end{block}
	\end{frame}
	
	%------------------------------------------------
	
	\begin{frame}{Cociente de potencias de igual base}
		\begin{idea}
			Cuando se dividen potencias de igual base, se conserva la base y se restan los exponentes.
		\end{idea}		
		\[
		\frac{a^m}{a^n}=a^{m-n},
		\qquad a\neq 0
		\]
		
		\pause
		
		\begin{exampleblock}{Ejemplo contextualizado}
			Un índice de atención se modela como: $\dfrac{x^8}{x^3}$
			Entonces: $\dfrac{x^8}{x^3}=x^{8-3}=x^5$
		\end{exampleblock}
		
		\pause
		
		\begin{alerta}
			La restricción es:$x\neq 0$, porque no existe división por cero.
		\end{alerta}
	\end{frame}
	
	%------------------------------------------------
	
	\begin{frame}{Potencia de una potencia}
		\begin{idea}
			Cuando una potencia se eleva a otra potencia, se multiplican los exponentes.
		\end{idea}
		
		\[
		(a^m)^n=a^{mn}
		\]
		
		\pause
		
		\begin{exampleblock}{Ejemplo en memoria de trabajo}
			Si un modelo de desempeño en memoria se expresa como:
			$(M^2)^3$
				
			entonces: $(M^2)^3=M^{2\cdot 3}=M^6$
		\end{exampleblock}
		
		\pause
		
		\begin{block}{Interpretación}
			La variable \(M\) representa un indicador de memoria de trabajo sometido a transformaciones sucesivas.
		\end{block}
	\end{frame}
	
	%------------------------------------------------
	
	\begin{frame}{Potencia de un producto}
		\begin{idea}
			La potencia de un producto se distribuye sobre cada factor.
		\end{idea}		
		\[
		(ab)^n=a^n b^n
		\]
		
		\pause
		
		\begin{exampleblock}{Ejemplo contextualizado}
			Suponga que \(a\) mide activación fisiológica y \(b\) mide demanda cognitiva. Entonces:			
			\[
			(2ab)^3=2^3a^3b^3=8a^3b^3
			\]
		\end{exampleblock}
		
		\pause
		
		\begin{alerta}
			No confundir: $(a+b)^2\neq a^2+b^2$. Esto es un error frecuente al analizar expresiones algebraicas.
		\end{alerta}
	\end{frame}
	
	%------------------------------------------------
	
	\begin{frame}{Exponente cero y negativo}
		\begin{block}{Exponente cero}
			\[
			a^0=1,\qquad a\neq 0
			\]
		\end{block}
		
		\pause
		
		\begin{exampleblock}{Ejemplo}
			\[
			E^0=1
			\]
			
			si \(E\neq 0\), donde \(E\) puede representar un indicador de estrés.
		\end{exampleblock}
		
		\pause
		
		\begin{block}{Exponente negativo}
			$a^{-n}=\dfrac{1}{a^n},	\qquad a\neq 0$
		\end{block}
		
		\pause
		
		\begin{exampleblock}{Ejemplo}
			$T^{-2}=\dfrac{1}{T^2}$, donde \(T\) puede representar tiempo de respuesta.
		\end{exampleblock}
	\end{frame}
	
	%------------------------------------------------
	\section{Errores e inconsistencias}
	
	\begin{frame}{Error frecuente 1}
		Un estudiante escribe:
		
		\[
		A^3+A^5=A^8
		\]
		
		\pause
		
		\begin{alerta}
			Esto es incorrecto.
		\end{alerta}
		
		\pause
		
		La propiedad:
		
		\[
		A^m\cdot A^n=A^{m+n}
		\]
		
		se aplica solamente en multiplicaciones, no en sumas.
		
	\end{frame}
	
	%------------------------------------------------
	
	\begin{frame}{Error frecuente 2}
		En un modelo de ansiedad, un estudiante plantea:
		
		\[
		(3A^2)^3=3A^6
		\]
		
		\pause
		
		\begin{alerta}
			El resultado es incorrecto.
		\end{alerta}
		
		\pause
		
		La potencia afecta tanto al coeficiente como a la variable:
		
		\[
		(3A^2)^3=3^3(A^2)^3
		\]
		
		\[
		=27A^6
		\]
		
		\pause
		
		\begin{block}{Resultado correcto}
			\[
			(3A^2)^3=27A^6
			\]
		\end{block}
	\end{frame}
	 %%%%%%%%%%%%%%%%%%%%%%%%%%%%%%%%%%%%%%%%%%%%%%%%%%%%%%%%%%%%%%%%
	
	\section{Potencias en contextos psicológicos}
	
	%------------------------------------------------
	\begin{frame}{Ejemplo 1: carga cognitiva creciente}
		En una tarea experimental, la carga cognitiva se modela por:
		\[
		C(n)=3\cdot 2^{n-1}
		\]
		donde \(n\) es el nivel de dificultad. Calcular $C(5)$.
		
		\[
		C(5)=3\cdot 2^{5-1}
		\]
		
		\[
		C(5)=3\cdot 2^4
		\]
		
		\[
		C(5)=3\cdot 16=48
		\]
		
		\begin{block}{Interpretación}
			En el nivel 5, la carga cognitiva relativa es de \(48\) unidades.
		\end{block}
	\end{frame}
	
	%------------------------------------------------
	\begin{frame}{Ejercicio 1: carga cognitiva}
		Una tarea de memoria visual comienza con \(4\) estímulos. En cada nivel, la cantidad de estímulos se triplica.
		
		\begin{enumerate}
			\item Modele la cantidad de estímulos en el nivel \(n\).
			\item Calcule la cantidad de estímulos en el nivel \(6\).
		\end{enumerate}
		
		\pause
		
		\begin{block}{Solución}
			\[
			E(n)=4\cdot 3^{n-1}
			\]
			
			\[
			E(6)=4\cdot 3^5=4\cdot 243=972
			\]
			
			El crecimiento es exponencial porque cada nivel multiplica por \(3\) al anterior.
		\end{block}
	\end{frame}
	
	%------------------------------------------------
	\begin{frame}{Ejemplo 2: reducción de ansiedad}
		Después de cada sesión terapéutica, un indicador de ansiedad se reduce al \(80\%\) del valor anterior.
		
		\[
		A(n)=A_0(0{,}8)^n
		\]
		
		Si \(A_0=75\), entonces después de \(5\) sesiones:
		
		\[
		A(5)=75(0{,}8)^5
		\]
		
		\[
		A(5)=75(0{,}32768)
		\]
		
		\[
		A(5)=24{,}576
		\]
		
		\begin{block}{Interpretación}
			Después de 5 sesiones, el indicador estimado de ansiedad es aproximadamente \(24{,}58\).
		\end{block}
	\end{frame}
	
	%------------------------------------------------
	
	%------------------------------------------------
	\begin{frame}{Ejemplo 3: memoria y olvido}
		Un modelo simple de retención de información (en porcentaje) está dado por:
		\[
		R(t)=100(0{,}7)^t
		\]
		donde \(t\) representa días desde el aprendizaje inicial.
		
		Determinar el porcentaje de retensión después de 3 días.
		
		\[
		R(3)=100(0{,}7)^3
		\]
		
		\[
		R(3)=100(0{,}343)=34{,}3
		\]
		
		\begin{block}{Interpretación}
			Después de 3 días, se retiene aproximadamente el \(34{,}3\%\) de la información inicial.
		\end{block}
	\end{frame}
	
	%------------------------------------------------
	\begin{frame}{Ejercicio 3: retención de información}
		Un estudiante retiene el \(92\%\) de lo aprendido cada día. Inicialmente domina el \(100\%\) de un contenido.
		
		\begin{enumerate}
			\item Plantee el modelo de retención.
			\item Calcule la retención después de \(7\) días.
			\item Interprete el resultado.
		\end{enumerate}
		
		\pause
		
		\begin{block}{Solución}
			\[
			\begin{array}{rcl}
				R(t)   & = & 100(0{,}92)^t \\[0.3cm]
				R(7)   & = & 100(0{,}92)^7 \\[0.3cm]
				R(7)   & \approx & 100(0{,}5578) \\[0.3cm]
				R(7)   & \approx & 55{,}78
			\end{array}
			\]
			
			Después de 7 días, retiene aproximadamente el \(55{,}78\%\).
		\end{block}
	\end{frame}
	
	%------------------------------------------------
	\begin{frame}{Ejemplo 4: simplificación algebraica con potencias}
		Un índice combinado de ansiedad \(A\), memoria \(M\) y atención \(T\) está dado por:
		
		\[
		I=\frac{12A^5M^3T^4}{3A^2MT}
		\]
		
		\[
		I=4A^{5-2}M^{3-1}T^{4-1}
		\]
		
		\[
		I=4A^3M^2T^3
		\]
		
		\begin{block}{Restricciones}
			\[
			A\neq 0,\qquad M\neq 0,\qquad T\neq 0
			\]
		\end{block}
	\end{frame}
	
	%------------------------------------------------
	\begin{frame}{Ejercicio 4: índice psicológico compuesto}
		Simplifique el siguiente índice:
		
		\[
		P=\frac{18E^7A^4M^5}{6E^2A^3M}
		\]
		
		donde \(E\) representa estrés, \(A\) ansiedad y \(M\) memoria de trabajo.
		
		\pause
		
		\begin{block}{Solución}
			\[
			P=3E^{7-2}A^{4-3}M^{5-1}
			\]
			
			\[
			P=3E^5AM^4
			\]
			
			\[
			E\neq 0,\qquad A\neq 0,\qquad M\neq 0
			\]
		\end{block}
	\end{frame}
	
	%------------------------------------------------
	\begin{frame}{Ejemplo 5: error frecuente en potencia de producto}
		Un estudiante simplifica:
		
		\[
		(4A^2M^3)^2
		\]
		
		como:
		
		\[
		4A^4M^6
		\]
		
		\begin{alertblock}{Error}
			No elevó el coeficiente \(4\) al cuadrado.
		\end{alertblock}
		
		\[
		(4A^2M^3)^2
		=
		4^2(A^2)^2(M^3)^2
		\]
		
		\[
		=16A^4M^6
		\]
	\end{frame}
	
	%------------------------------------------------
	\begin{frame}{Ejercicio 5: detectar inconsistencia}
		Un estudiante afirma que:
		
		\[
		(5E^3T^2)^2=5E^6T^4
		\]
		
		\begin{enumerate}
			\item Identifique el error.
			\item Corrija el desarrollo.
			\item Interprete por qué el coeficiente también debe elevarse.
		\end{enumerate}
		
		\pause
		
		\begin{block}{Solución}
			\[
			(5E^3T^2)^2
			=
			5^2(E^3)^2(T^2)^2
			\]
			
			\[
			=25E^6T^4
			\]
			
			El coeficiente forma parte del producto completo.
		\end{block}
	\end{frame}
	
	%------------------------------------------------
	\begin{frame}{Ejercicio 6: patrón exponencial en difusión social}
		En una simulación de influencia social, cada participante comparte una recomendación con \(3\) nuevas personas. Inicialmente participan \(6\) personas. El modelo del número de personas alcanzadas en la etapa $n$ es 
			\[
		P(n)=6\cdot 3^{n-1}
		\]
	Calcule el número de personas alcanzadas en la etapa \(5\).
					
		\pause
		
		\begin{block}{Solución}
			\[
			P(5)=6\cdot 3^4=6\cdot81=486
			\]
			
		\end{block}
	\end{frame}
	
	%------------------------------------------------
	
	%------------------------------------------------
	\begin{frame}{Ejercicio 7: comparación de modelos}
		Dos programas de intervención reducen un indicador de estrés de la siguiente forma:
		\[
		E_1(n)=80(0{,}9)^n
		\]
		
		\[
		E_2(n)=80(0{,}75)^n
		\]
		
		\begin{enumerate}
			\item Calcule ambos valores para \(n=4\).
			\item Determine cuál reduce más rápido el estrés.
			\item Explique usando potencias.
		\end{enumerate}
		
		\pause
		
		\begin{block}{Solución}
			\[
			E_1(4)=80(0{,}9)^4=80(0{,}6561)=52{,}488
			\]
			
			\[
			E_2(4)=80(0{,}75)^4=80(0{,}31640625)=25{,}3125
			\]
			
			El segundo programa reduce más rápido el estrés.
		\end{block}
	\end{frame}
	
	%------------------------------------------------
	\begin{frame}{Ejercicio 8}
		Simplifique:
		
		\[
		K=
		\frac{(6A^3M^2)^2}{3A^2M}
		\]
		
		\pause
		
		Primero:
		
		\[
		(6A^3M^2)^2
		=
		36A^6M^4
		\]
		
		Luego:
		
		\[
		K=\frac{36A^6M^4}{3A^2M}
		\]
		
		\[
		K=12A^{6-2}M^{4-1}
		\]
		
		\[
		K=12A^4M^3
		\]
		
		
	\end{frame}
	
	%------------------------------------------------
	\begin{frame}{Ejercicio 9: modelo integrado}
		Un índice experimental se define por: $	I=\dfrac{(2E^2A^3)^3}{4EA^2}$
		
		donde \(E\) representa estrés y \(A\) ansiedad.
		
		\begin{enumerate}
			\item Simplifique el índice.
			\item Determine restricciones.
			\item Explique qué propiedad de potencias se utiliza en cada paso.
		\end{enumerate}
		
		\pause
		
		\begin{block}{Solución}
			\[
			\begin{array}{rcl}
				(2E^2A^3)^3 & = & 2^3E^6A^9 \\[0.3cm]
				& = & 8E^6A^9 \\[0.3cm]
				I           & = & \dfrac{8E^6A^9}{4EA^2} \\[0.3cm]
				I           & = & 2E^5A^7
			\end{array}
			\]
			
			Restricciones:$	E\neq0,\quad A\neq0$
		\end{block}
	\end{frame}
	
	%%%%%%%%%%%%%%%%%%%%%%%%%%%%%%%%%%%%%%%%%%%%%%%%%%%%%%%%%%%%%%%
	%------------------------------------------------
	\section{Elementos básicos del álgebra}
		
	\begin{frame}{Término algebraico}
		\begin{idea}
			Un término algebraico es una expresión formada por números y letras unidos mediante multiplicación y potencias.
		\end{idea}
		
		\[
		-5x^3y^2
		\]
		
		\pause
		
		\begin{block}{Elementos}
			\begin{itemize}
				\item Coeficiente numérico: \(-5\).
				\item Parte literal: \(x^3y^2\).
				\item Variables: \(x\), \(y\).
				\item Exponentes: \(3\), \(2\).
			\end{itemize}
		\end{block}
		
		\pause
		
	\end{frame}
	
	%------------------------------------------------
	\section{Elementos básicos del álgebra}
	\begin{frame}{Expresión algebraica}
		\begin{idea}
			Una expresión algebraica combina términos algebraicos mediante operaciones.
		\end{idea}
		\[
		4x^2-7x+3
		\]
		
		\pause
		
		\begin{exampleblock}{Ejemplo contextualizado}
			Un puntaje estimado de desempeño puede modelarse mediante:
					\[
			P(t)=4t^2-7t+3
			\]
			donde \(t\) representa el tiempo de práctica en semanas.
		\end{exampleblock}
	\end{frame}	
	
	\begin{frame}{Diferencia ente expresión algebraica y ecuación}
		
		\begin{alerta}
			Una expresión algebraica no es necesariamente una ecuación.
			
			\[
			4t^2-7t+3
			\]
			
			es una expresión.
			
			\[
			4t^2-7t+3=20
			\]
			
			es una ecuación.
		\end{alerta}
	\end{frame}
	
	%------------------------------------------------
	
	\begin{frame}{Términos semejantes}
		\begin{idea}
			Dos términos son semejantes si tienen exactamente la misma parte literal.
		\end{idea}
		
		\[
		7x^2y
		\qquad\text{y}\qquad
		-3x^2y
		\]
		
		son semejantes.
		
		\pause
		
		\[
		5xy^2
		\qquad\text{y}\qquad
		5x^2y
		\]
		
		no son semejantes.
		
		\pause
		
		\begin{exampleblock}{Contexto psicológico}
			Si \(A\) representa ansiedad y \(M\) motivación:
			
			\[
			8A^2M
			\qquad\text{y}\qquad
			-3A^2M \qquad 
			 \text{son términos semejantes}
			 \]
		\end{exampleblock}
	\end{frame}
	
	%------------------------------------------------
	
	\begin{frame}{Reducción de términos semejantes}
		\begin{exampleblock}{Ejemplo}
			Simplifique el siguiente modelo de puntaje psicológico:
			
			\[
			P=8A^2-3A+5A^2+9A-4
			\]
			
			donde \(A\) representa un indicador de ansiedad.
		\end{exampleblock}
		
		\pause
		
		Agrupamos términos semejantes:
		
		\[
		P=(8A^2+5A^2)+(-3A+9A)-4
		\]
		
		\pause
		
		\[
		P=13A^2+6A-4
		\]
		
		\pause
		
		\begin{block}{Resultado}
			\[
			P=13A^2+6A-4
			\]
		\end{block}
	\end{frame}
	
	%------------------------------------------------
	\section{Patrones y relaciones matemáticas}
	
	\begin{frame}{Patrón lineal: sesiones terapéuticas}
		Un estudiante registra el número acumulado de ejercicios de autorregulación realizados por un paciente:
		
		\[
		\begin{array}{c|c}
			\text{Semana} & \text{Ejercicios acumulados} \\
			\hline
			1 & 5 \\
			2 & 8 \\
			3 & 11 \\
			4 & 14 \\
			5 & 17
		\end{array}
		\]
		
		\pause
		
		\begin{block}{Patrón}
			El número aumenta de \(3\) en \(3\). \qquad  $a_n=3n+2$
		\end{block}
		
		\pause
		
		\begin{block}{Interpretación}
			En la semana \(n\), el paciente acumula \(3n+2\) ejercicios.
		\end{block}
	\end{frame}
	
	%------------------------------------------------
	
	\begin{frame}{Patrón cuadrático: tareas cognitivas}
		En una prueba de memoria, la cantidad de combinaciones posibles aumenta como:
		
		\[
		1,\ 4,\ 9,\ 16,\ 25,\ldots
		\]
		
		\pause
		
		\[
		1^2,\ 2^2,\ 3^2,\ 4^2,\ 5^2,\ldots
		\]
		
		\pause
		
		\begin{block}{Modelo}
			\[
			a_n=n^2
			\]
		\end{block}
		
		\pause
		
		\begin{actividad}
			En una tarea de asociación de palabras, el número de relaciones posibles se modela por:
			
			\[
			a_n=2n^2
			\]
			
			Calcule \(a_3\), \(a_4\) y \(a_5\).
		\end{actividad}
	\end{frame}
	
	%------------------------------------------------
	
	\begin{frame}{Patrón exponencial: propagación de respuestas}
		En una simulación de dinámica grupal, cada participante transmite una instrucción a dos nuevos participantes.
		
		\[
		5,\ 10,\ 20,\ 40,\ 80,\ldots
		\]
		
		\pause
		
		\[
		5,\quad 5\cdot2,\quad 5\cdot2^2,\quad 5\cdot2^3,\quad 5\cdot2^4
		\]
		
		\pause
		
		\begin{block}{Modelo}
			\[
			a_n=5\cdot2^{n-1}
			\]
		\end{block}
		
		\pause
		
		\begin{alerta}
			El exponente es \(n-1\), porque el primer término corresponde a:
			
			\[
			5\cdot2^0=5
			\]
		\end{alerta}
	\end{frame}
	
	%------------------------------------------------
	\section{Ejemplos contextualizados}
	
	\begin{frame}{Ejemplo 1: registro de respuestas}
		En un experimento, el número de respuestas registradas aumenta al doble en cada bloque.  
		Inicialmente se registran \(120\) respuestas.
		
		\pause
		
		\begin{block}{Modelo}
			\[
			R(t)=120\cdot2^t
			\]
		\end{block}
		
		\pause
		
		\begin{exampleblock}{Pregunta}
			¿Cuántas respuestas se esperan después de \(5\) bloques?
		\end{exampleblock}
		
		\pause
		
	\[
	\begin{array}{rcl}
		R(5) & = & 120\cdot 2^5 \\[0.3cm]
		& = & 120\cdot 32 \\[0.3cm]
		& = & 3840
	\end{array}
	\]
		
		\pause
		
		\textbf{Respuesta:} 
			Después de \(5\) bloques se esperan \(3840\) respuestas registradas.
	\end{frame}
	
	%------------------------------------------------
	
	\begin{frame}{Ejemplo 2: índice combinado de atención}
		Un índice experimental combina atención sostenida \(A\) y memoria de trabajo \(M\):
		
		\[
		I=4A^3M^2\cdot3A^2M^5
		\]
		
		\pause
		
		Multiplicamos coeficientes:
		
		\[
		4\cdot3=12
		\]
		
		\pause
		
		Aplicamos propiedades de potencias:
		
		\[
		A^3\cdot A^2=A^5
		\]
		
		\[
		M^2\cdot M^5=M^7
		\]
		
		\pause
		
		\begin{block}{Resultado}
			\[
			I=12A^5M^7
			\]
		\end{block}
	\end{frame}
	
	%------------------------------------------------
	
	\begin{frame}{Ejemplo 3: razón entre indicadores}
		Un modelo compara dos indicadores de respuesta emocional:
		
		\[
		Q=\frac{18E^7T^4}{6E^3T}
		\]
		
		donde \(E\) representa intensidad emocional y \(T\) tiempo de respuesta.
		
		\pause
		
		\[
		\begin{array}{rcl}
			\dfrac{18}{6} & = & 3 \\[0.3cm]
			\dfrac{E^7}{E^3} & = & E^4 \\[0.3cm]
			\dfrac{T^4}{T} & = & T^3
		\end{array}
		\]
		
		\pause
		
		\begin{block}{Resultado}
			\[
			Q=3E^4T^3
			\]
		\end{block}
		
		\pause
		
		\begin{alerta}
			Restricciones:
			
			\[
			E\neq0,\qquad T\neq0
			\]
		\end{alerta}
	\end{frame}
	
	%------------------------------------------------
	
	\begin{frame}{Ejemplo 4: escala compuesta}
		Una escala simplificada de ajuste psicológico se representa por:
		
		\[
		S=7A^2M-3AM^2+5A^2M+9AM^2-4M
		\]
		
		donde:
		
		\begin{itemize}
			\item \(A\): ansiedad.
			\item \(M\): motivación.
		\end{itemize}
		
		\pause
		
		Agrupamos términos semejantes:
		
		\[
		S=(7A^2M+5A^2M)+(-3AM^2+9AM^2)-4M
		\]
		
		\pause
		
		\[
		S=12A^2M+6AM^2-4M
		\]
		
		\pause
		
		\begin{block}{Resultado}
			\[
			S=12A^2M+6AM^2-4M
			\]
		\end{block}
	\end{frame}
	
	%------------------------------------------------
	\section{Actividades}
	
	\begin{frame}{Actividad 1: potencias en contexto psicológico}
		Simplifique las siguientes expresiones:
		
		\[
		\begin{array}{ll}
			\text{a)} & A^4\cdot A^7 \\[0.2cm]
			\text{b)} & \dfrac{M^{10}}{M^3} \\[0.3cm]
			\text{c)} & (E^5)^4 \\[0.2cm]
			\text{d)} & (2A^3M)^2 \\[0.2cm]
			\text{e)} & \dfrac{12E^6T^4}{3E^2T}
		\end{array}
		\]
		
		\pause
		
		\begin{block}{Respuestas}
			\[
			\begin{array}{ll}
				\text{a)} & A^{11} \\[0.2cm]
				\text{b)} & M^7 \\[0.2cm]
				\text{c)} & E^{20} \\[0.2cm]
				\text{d)} & 4A^6M^2 \\[0.2cm]
				\text{e)} & 4E^4T^3
			\end{array}
			\]
		\end{block}
	\end{frame}
	
	%------------------------------------------------
	
	\begin{frame}{Actividad 2: reducción algebraica}
		Un modelo simplificado de desempeño académico-emocional está dado por:
		
		\[
		D=9A^2-5A+4-3A^2+11A-8
		\]
		
		\pause
		
		Agrupamos términos semejantes:
		
		\[
		D=(9A^2-3A^2)+(-5A+11A)+(4-8)
		\]
		
		\pause
		
		\[
		D=6A^2+6A-4
		\]
		
		\pause
		
		\begin{block}{Resultado}
			\[
			D=6A^2+6A-4
			\]
		\end{block}
	\end{frame}
	
	%------------------------------------------------
	
	\begin{frame}{Actividad 3: detectar inconsistencias}
		Un estudiante simplifica el modelo:
		
		\[
		(4M^2)^3
		\]
		
		de la siguiente manera:
		
		\[
		(4M^2)^3=4M^6
		\]
		
		\pause
		
		\begin{alerta}
			Detecte el error.
		\end{alerta}
		
		\pause
		
		La potencia afecta al coeficiente y a la variable:
		
		\[
		(4M^2)^3=4^3(M^2)^3
		\]
		
		\[
		=64M^6
		\]
		
		\pause
		
		\begin{block}{Resultado correcto}
			\[
			(4M^2)^3=64M^6
			\]
		\end{block}
	\end{frame}
	
	%------------------------------------------------
	
	\begin{frame}{Actividad 4: patrón en intervención psicológica}
		En un programa de intervención, el número de registros conductuales sigue la secuencia:
		
		\[
		4,\ 12,\ 36,\ 108,\ 324,\ldots
		\]
		
		\pause
		
		\begin{block}{Pregunta}
			Determine el patrón y proponga el término general.
		\end{block}
		
		\pause
		
		Cada término se multiplica por \(3\):
		
		\[
		4,\quad 4\cdot3,\quad 4\cdot3^2,\quad 4\cdot3^3
		\]
		
		\pause
		
		\[
		a_n=4\cdot3^{n-1}
		\]
		
		\pause
		
		\begin{block}{Interpretación}
			En la sesión \(n\), el número estimado de registros es:
			
			\[
			4\cdot3^{n-1}
			\]
		\end{block}
	\end{frame}
	
	%------------------------------------------------
	
	\begin{frame}{Actividad 5: caso integrador}
		Un equipo de investigación psicológica propone el siguiente índice:
		
		\[
		I=\frac{6A^4M^3}{2A^2M}\,+\,5A^2M^2-3A^2M^2
		\]
		
		donde:
		
		\begin{itemize}
			\item \(A\): ansiedad.
			\item \(M\): memoria de trabajo.
		\end{itemize}
		
		\pause
		
		Simplificamos el cociente:
		
		\[
		\frac{6A^4M^3}{2A^2M}
		=
		3A^2M^2
		\]
		
		\pause
		
		Entonces:
		
		\[
		I=3A^2M^2+5A^2M^2-3A^2M^2
		\]
		
		\pause
		
		\[
		I=5A^2M^2
		\]
		
		\pause
		
		\begin{block}{Resultado}
			\[
			I=5A^2M^2
			\]
		\end{block}
	\end{frame}
	
	%------------------------------------------------
	\section{Síntesis final}
	
	\begin{frame}{Síntesis conceptual}
		\begin{block}{Potencias}
			Permiten representar multiplicaciones reiteradas y modelar crecimiento, escalamiento y patrones.
		\end{block}
		
		\pause
		
		\begin{block}{Álgebra}
			Un término algebraico posee coeficiente, variables y exponentes.  
			Una expresión algebraica combina términos mediante operaciones.
		\end{block}
		
		\pause
		
		\begin{block}{Psicología}
			Las expresiones algebraicas pueden representar modelos simplificados de atención, ansiedad, memoria, motivación o desempeño.
		\end{block}
		
		\pause
		
		\begin{block}{Pensamiento crítico}
			Detectar errores algebraicos permite interpretar mejor los modelos y evitar conclusiones incorrectas.
		\end{block}
	\end{frame}
	
	%------------------------------------------------
	
	\begin{frame}{Cierre}
		\begin{center}
			{\Huge \textbf{Conclusión}}
			
			\vspace{0.5cm}
			
			\Large
			Las potencias y el álgebra permiten representar patrones, simplificar modelos y analizar relaciones presentes en fenómenos psicológicos.
			
			\vspace{0.7cm}
			
			\textbf{Pregunta final:}
			
			\[
			A^2+A^2=2A^2
			\]
			
			pero
			
			\[
			A^2\cdot A^2=A^4
			\]
			
			\vspace{0.4cm}
			
			¿Por qué ambas operaciones producen resultados distintos?
		\end{center}
	\end{frame}
	
	%------------------------------------------------
\end{document}